\documentclass[11pt,a4paper]{article}

% 한글 지원 (XeLaTeX)
\usepackage{fontspec}
\usepackage{kotex}
\setmainfont{Noto Sans CJK KR}
\setsansfont{Noto Sans CJK KR}
\setmonofont{Noto Sans Mono CJK KR}

% 패키지
\usepackage{graphicx}
\usepackage{xcolor}
\usepackage{booktabs}
\usepackage{tabularx}
\usepackage{longtable}
\usepackage{float}
\usepackage{caption}
\usepackage{hyperref}
\usepackage{geometry}
\usepackage{fancyhdr}
\usepackage{tcolorbox}
\usepackage{enumitem}

% 페이지 설정
\geometry{margin=2.5cm}
\pagestyle{fancy}
\fancyhf{}
\fancyhead[L]{사과 가격 예측 시뮬레이터 사용설명서}
\fancyhead[R]{v1.2.1.1}
\fancyfoot[C]{\thepage}

% 색상 정의
\definecolor{appleRed}{HTML}{E53935}
\definecolor{appleGreen}{HTML}{43A047}
\definecolor{honeyYellow}{HTML}{FFC107}
\definecolor{tipBg}{HTML}{FFFDE7}
\definecolor{tipBorder}{HTML}{FFC107}
\definecolor{grayBg}{HTML}{F5F5F5}

% 팁 박스
\newtcolorbox{tipbox}[1][]{
  colback=tipBg,
  colframe=tipBorder,
  fonttitle=\bfseries,
  title={💡 Tip},
  #1
}

% 하이퍼링크 설정
\hypersetup{
  colorlinks=true,
  linkcolor=appleRed,
  urlcolor=appleRed
}

% 캡션 설정
\captionsetup{font=small, labelfont=bf}
\captionsetup[figure]{name=그림}
\captionsetup[table]{name=표}

\begin{document}

% 제목 페이지
\begin{titlepage}
\centering
\vspace*{3cm}
{\Huge\bfseries 🍎 사과 가격 예측 시뮬레이터\\[0.5cm]
사용설명서}\\[2cm]
{\Large\color{appleRed} v1.2.1.1}\\[1cm]
{\large 선형 회귀 학습 시뮬레이터}\\[3cm]
{\large 2026년 2월 1일}\\[2cm]
\vfill
{\small 청주교육대학교 컴퓨팅사고력 및 AI 교육}
\end{titlepage}

\tableofcontents
\newpage

%========================================
\section{📘 가이드 사용법}
%========================================

이 설명서를 효과적으로 활용하는 방법입니다.

\begin{table}[H]
\centering
\caption{가이드 사용 순서}
\label{tab:guide_usage}
\begin{tabularx}{\textwidth}{c|l|X}
\toprule
\textbf{순서} & \textbf{섹션} & \textbf{내용} \\
\midrule
1단계 & 제1장 개요 & 시뮬레이터의 목적과 전체 구조 파악 \\
2단계 & 제3장 품질 점수 & 핵심 개념인 품질 점수 이해 \\
3단계 & 제4장 기능별 사용법 & 실제 조작 방법 학습 \\
4단계 & 제6장 학습 결과 해석 & 결과 분석 방법 이해 \\
\bottomrule
\end{tabularx}
\end{table}

\begin{tipbox}
처음 사용하시는 분은 \textbf{1장 → 3장 → 4장} 순서로 읽으시면 됩니다. 
용어가 어려우면 \textbf{제8장 용어 해설}을 먼저 보세요!
\end{tipbox}

%========================================
\section{개요}
%========================================

\subsection{시뮬레이터 소개}

\textbf{사과 가격 예측 시뮬레이터}는 머신러닝의 핵심 알고리즘인 \textbf{선형 회귀(Linear Regression)}의 
원리를 직관적으로 이해할 수 있도록 설계된 교육용 도구입니다.

\subsection{교육 목표}

\begin{table}[H]
\centering
\caption{교육 목표}
\label{tab:edu_goals}
\begin{tabularx}{\textwidth}{l|X}
\toprule
\textbf{목표} & \textbf{설명} \\
\midrule
선형 회귀 이해 & 입력 변수와 출력 변수 간의 선형 관계 학습 \\
경사 하강법 체험 & 손실 함수를 최소화하는 최적화 과정 시각화 \\
예측 원리 학습 & 학습된 모델로 새로운 데이터 예측 체험 \\
\bottomrule
\end{tabularx}
\end{table}

\subsection{시스템 요구사항}

\begin{table}[H]
\centering
\caption{시스템 요구사항}
\label{tab:system_req}
\begin{tabularx}{\textwidth}{l|X}
\toprule
\textbf{항목} & \textbf{요구사항} \\
\midrule
웹 브라우저 & Chrome, Edge, Firefox (최신 버전 권장) \\
화면 해상도 & 1280 × 720 이상 \\
인터넷 연결 & 초기 로딩 시 필요 (CDN 라이브러리) \\
\bottomrule
\end{tabularx}
\end{table}

%========================================
\section{화면 구성}
%========================================

\subsection{모드 전환}

시뮬레이터는 \textbf{학습 모드}와 \textbf{추론 모드} 두 가지로 구성됩니다.

\begin{figure}[H]
\centering
\includegraphics[width=0.7\textwidth]{fig_modeswitch.png}
\caption{모드 전환 버튼}
\label{fig:modeswitch}
\end{figure}

\begin{itemize}
  \item \textbf{🧠 학습 모드}: AI가 데이터에서 패턴을 찾아 학습하는 단계
  \item \textbf{🔮 추론 모드}: 학습 완료 후 새로운 값으로 예측하는 단계
\end{itemize}

\begin{tipbox}
\textbf{학습}은 ``공부하는 것'', \textbf{추론}은 ``배운 걸로 문제 푸는 것''이라고 생각하세요!
\end{tipbox}

\subsection{학습 모드 화면}

\begin{figure}[H]
\centering
\includegraphics[width=0.9\textwidth]{fig_trainmode.png}
\caption{학습 모드 전체 화면}
\label{fig:trainmode}
\end{figure}

학습 모드는 3열 레이아웃으로 구성됩니다:

\begin{table}[H]
\centering
\caption{학습 모드 화면 구성}
\label{tab:trainmode_layout}
\begin{tabularx}{\textwidth}{l|l|X}
\toprule
\textbf{위치} & \textbf{패널} & \textbf{기능} \\
\midrule
왼쪽 & 📊 데이터 생성 & 데이터 개수, 노이즈 설정 \\
중앙 & 📈 차트 영역 & 손실 그래프, 산점도 \\
오른쪽 & ⚙️ 학습 설정 & 학습률, 에포크 조절 \\
\bottomrule
\end{tabularx}
\end{table}

\subsection{추론 모드 화면}

\begin{figure}[H]
\centering
\includegraphics[width=0.9\textwidth]{fig_infermode.png}
\caption{추론 모드 전체 화면}
\label{fig:infermode}
\end{figure}

\begin{table}[H]
\centering
\caption{추론 모드 화면 구성}
\label{tab:infermode_layout}
\begin{tabularx}{\textwidth}{l|l|X}
\toprule
\textbf{위치} & \textbf{패널} & \textbf{기능} \\
\midrule
왼쪽 & 🎯 예측 입력 & 당도, 무게 슬라이더 \\
중앙 & 📈 예측 시각화 & 산점도 + 회귀선 + 예측점 \\
오른쪽 & 💰 예측 결과 & 계산된 가격, 추론 과정 \\
\bottomrule
\end{tabularx}
\end{table}

%========================================
\section{품질 점수 이해하기}
%========================================

\subsection{품질 점수란?}

이 시뮬레이터의 핵심 개념입니다. \textbf{당도}와 \textbf{무게}, 두 가지 특성을 
하나의 \textbf{품질 점수}로 통합합니다.

\begin{tipbox}
\textbf{왜 품질 점수로 통합할까요?}\\[0.3cm]
두 변수(당도, 무게)를 따로 처리하면 3D 그래프가 필요해서 이해하기 어렵습니다.
하나의 점수로 합치면 \textbf{2D 그래프}로 쉽게 볼 수 있습니다!
\end{tipbox}

\subsection{계산 공식}

\begin{tcolorbox}[colback=grayBg, colframe=gray, title=품질 점수 계산식]
\textbf{품질 = 당도점수(50\%) + 무게점수(50\%)}\\[0.3cm]
\begin{itemize}[leftmargin=*]
  \item 당도점수 = (당도 - 10) ÷ 8 × 50 \quad (10\textasciitilde18 Brix → 0\textasciitilde50점)
  \item 무게점수 = (무게 - 150) ÷ 200 × 50 \quad (150\textasciitilde350g → 0\textasciitilde50점)
\end{itemize}
\end{tcolorbox}

\subsection{계산 예시}

\begin{table}[H]
\centering
\caption{품질 점수 계산 예시}
\label{tab:quality_example}
\begin{tabular}{l|c|c}
\toprule
\textbf{항목} & \textbf{값} & \textbf{점수} \\
\midrule
당도 & 14 Brix & (14-10)÷8×50 = 25점 \\
무게 & 250g & (250-150)÷200×50 = 25점 \\
\midrule
\textbf{총 품질} & & \textbf{50점} \\
\bottomrule
\end{tabular}
\end{table}

%========================================
\section{기능별 사용법}
%========================================

\subsection{데이터 생성}

\begin{figure}[H]
\centering
\includegraphics[width=0.4\textwidth]{fig_datagen.png}
\caption{데이터 생성 패널}
\label{fig:datagen}
\end{figure}

\begin{table}[H]
\centering
\caption{데이터 생성 옵션}
\label{tab:datagen_options}
\begin{tabularx}{\textwidth}{l|c|X}
\toprule
\textbf{항목} & \textbf{범위} & \textbf{설명} \\
\midrule
데이터 개수 & 5 \textasciitilde{} 100개 & 학습에 사용할 사과 데이터 수 \\
노이즈 수준 & 5\% / 10\% / 20\% & 현실의 불확실성 표현 \\
\bottomrule
\end{tabularx}
\end{table}

\begin{tipbox}
\textbf{노이즈}란 현실 데이터의 ``흔들림''입니다. 같은 품질이어도 가격이 조금씩 다를 수 있죠!
\end{tipbox}

\subsection{학습 실행}

\begin{figure}[H]
\centering
\includegraphics[width=0.9\textwidth]{fig_training.png}
\caption{학습 실행 화면}
\label{fig:training}
\end{figure}

\begin{table}[H]
\centering
\caption{학습 파라미터}
\label{tab:train_params}
\begin{tabularx}{\textwidth}{l|c|c|X}
\toprule
\textbf{파라미터} & \textbf{범위} & \textbf{권장값} & \textbf{설명} \\
\midrule
학습률 & 0.001 \textasciitilde{} 0.1 & 0.01 & 한 번에 얼마나 조정할지 \\
에포크 & 100 \textasciitilde{} 2000 & 500 & 전체 데이터 학습 반복 횟수 \\
\bottomrule
\end{tabularx}
\end{table}

\begin{tipbox}
\textbf{학습률}이 너무 크면 → 목표를 지나쳐 발산!\\
\textbf{학습률}이 너무 작으면 → 학습이 너무 느림
\end{tipbox}

\subsection{오차선 분석}

\begin{figure}[H]
\centering
\includegraphics[width=0.9\textwidth]{fig_errorline.png}
\caption{오차선(잔차) 표시 화면}
\label{fig:errorline}
\end{figure}

학습 완료 후 ``오차선 표시'' 체크박스를 활성화하면:

\begin{itemize}
  \item 실제 데이터 점과 예측선 사이의 거리(오차)를 시각화
  \item 녹색(작은 오차) → 빨간색(큰 오차) 색상 변화
  \item 평균 오차, 최대 오차, MSE 통계 표시
\end{itemize}

\begin{tipbox}
오차선을 통해 \textbf{``AI는 항상 정확하지 않다''}는 것을 눈으로 확인할 수 있습니다!
\end{tipbox}

\subsection{추론 (예측)}

\begin{figure}[H]
\centering
\includegraphics[width=0.9\textwidth]{fig_inference.png}
\caption{추론 모드 시각화}
\label{fig:inference}
\end{figure}

추론 모드에서는:

\begin{enumerate}
  \item 당도, 무게 슬라이더를 조절합니다
  \item 품질 점수가 자동으로 계산됩니다
  \item 예측 가격이 즉시 표시됩니다
  \item 차트에 노란 예측점이 표시됩니다
\end{enumerate}

%========================================
\section{시각화 요소}
%========================================

\subsection{산점도와 회귀선}

\begin{table}[H]
\centering
\caption{차트 요소 설명}
\label{tab:chart_elements}
\begin{tabularx}{\textwidth}{l|l|X}
\toprule
\textbf{요소} & \textbf{표시} & \textbf{의미} \\
\midrule
X축 & 품질 점수 & 0 \textasciitilde{} 100점 \\
Y축 & 가격 & 원 단위 \\
데이터 점 & 작은 원 & 실제 사과 데이터 (녹색→빨강: 품질) \\
회귀선 & 빨간 직선 & AI가 학습한 예측선 \\
예측점 & 노란 큰 원 & 현재 입력값의 예측 위치 \\
\bottomrule
\end{tabularx}
\end{table}

\subsection{손실 그래프}

\begin{itemize}
  \item \textbf{아래로 내려가면}: 학습이 잘 진행되는 중
  \item \textbf{평평해지면}: 학습 완료 (더 이상 개선 없음)
  \item \textbf{올라가면}: 학습률이 너무 큼! 낮춰보세요
\end{itemize}

%========================================
\section{학습 결과 해석}
%========================================

\subsection{회귀식 이해}

학습이 완료되면 다음과 같은 회귀식이 표시됩니다:

\begin{tcolorbox}[colback=grayBg, colframe=gray]
\centering
\textbf{가격 = 30.5 × 품질 + 985}
\end{tcolorbox}

이 수식의 의미:
\begin{itemize}
  \item 품질 1점 상승 → 가격 약 30원 상승
  \item 품질 0점일 때 기본 가격 약 985원
  \item 품질 100점이면 → 약 4,035원 예측
\end{itemize}

\subsection{R² 점수}

\begin{table}[H]
\centering
\caption{R² 점수 해석 기준}
\label{tab:r2_interpretation}
\begin{tabular}{c|l|l}
\toprule
\textbf{R² 범위} & \textbf{상태} & \textbf{의미} \\
\midrule
0.8 이상 & ✅ 양호 & 모델이 데이터를 잘 설명함 \\
0.5 \textasciitilde{} 0.8 & ⚠️ 보통 & 어느 정도 설명하지만 개선 필요 \\
0.5 미만 & ❌ 부족 & 데이터 추가 또는 설정 조정 필요 \\
\bottomrule
\end{tabular}
\end{table}

%========================================
\section{문제 해결 (FAQ)}
%========================================

\subsection{손실 그래프가 올라갔다 내려갔다 해요}

\begin{tcolorbox}[colback=tipBg, colframe=appleRed, title=Q: 손실 그래프가 불안정해요]
\textbf{A:} 학습률이 너무 큽니다!\\[0.2cm]
\textbf{해결:} 학습률을 0.01 이하로 낮춰보세요.
\end{tcolorbox}

\subsection{R² 점수가 음수예요}

\begin{tcolorbox}[colback=tipBg, colframe=appleRed, title=Q: R² 점수가 음수로 나와요]
\textbf{A:} 모델이 평균보다 못한 예측을 하고 있습니다.\\[0.2cm]
\textbf{해결:} 데이터를 더 생성하거나, 에포크를 늘려보세요.
\end{tcolorbox}

\subsection{학습이 너무 오래 걸려요}

\begin{tcolorbox}[colback=tipBg, colframe=appleRed, title=Q: 학습 시간이 너무 길어요]
\textbf{A:} 에포크가 너무 많거나 데이터가 많습니다.\\[0.2cm]
\textbf{해결:} 에포크를 500 이하로 줄여보세요.
\end{tcolorbox}

\subsection{예측 가격이 이상해요}

\begin{tcolorbox}[colback=tipBg, colframe=appleRed, title=Q: 예측 가격이 비정상적이에요]
\textbf{A:} 학습이 제대로 되지 않았을 수 있습니다.\\[0.2cm]
\textbf{해결:} 🗑️ 초기화 버튼을 누르고 다시 학습해보세요.
\end{tcolorbox}

%========================================
\section{용어 해설}
%========================================

\begin{longtable}{l|p{10cm}}
\caption{주요 용어 해설}
\label{tab:glossary}\\
\toprule
\textbf{용어} & \textbf{설명} \\
\midrule
\endfirsthead
\toprule
\textbf{용어} & \textbf{설명} \\
\midrule
\endhead
\bottomrule
\endfoot

선형 회귀 & 데이터의 패턴을 직선으로 표현하는 방법. 
\textit{비유: 점들 사이로 자를 대고 가장 잘 맞는 선 긋기} \\
\midrule

품질 점수 & 당도와 무게를 하나로 합친 종합 점수 (0\textasciitilde100점).
\textit{비유: 국어+수학 점수를 평균내서 ``총점''으로 보는 것} \\
\midrule

가중치 (w) & 품질이 가격에 미치는 영향력. 회귀선의 기울기.
\textit{비유: ``품질 1점당 가격 몇 원?''} \\
\midrule

편향 (b) & 품질이 0일 때의 기본 가격. 회귀선의 y절편.
\textit{비유: 기본 배송비처럼 항상 붙는 고정값} \\
\midrule

학습률 & 한 번의 학습에서 가중치를 얼마나 조정할지.
\textit{비유: 걸음 크기. 너무 크면 넘어지고, 너무 작으면 느림} \\
\midrule

에포크 & 전체 데이터를 한 번 학습하는 단위.
\textit{비유: 문제집을 처음부터 끝까지 1회 푸는 것} \\
\midrule

손실 (MSE) & 예측과 실제의 차이를 숫자로 표현. 작을수록 좋음.
\textit{비유: 과녁 맞추기에서 ``빗나간 거리''} \\
\midrule

R² 점수 & 모델이 데이터를 얼마나 잘 설명하는지 (0\textasciitilde1).
\textit{비유: 시험 점수. 1에 가까울수록 100점} \\
\midrule

오차 (잔차) & 실제값과 예측값의 차이.
\textit{비유: 예상 도착시간과 실제 도착시간의 차이} \\
\midrule

노이즈 & 데이터의 불규칙한 변동.
\textit{비유: 같은 품질 사과도 날마다 가격이 조금씩 다른 것} \\
\midrule

학습 모드 & AI가 데이터에서 패턴을 찾는 단계.
\textit{비유: 학생이 교과서로 공부하는 시간} \\
\midrule

추론 모드 & 학습된 모델로 새 값을 예측하는 단계.
\textit{비유: 공부 끝나고 시험 보는 시간} \\

\end{longtable}

%========================================
\section{참고자료}
%========================================

\begin{table}[H]
\centering
\caption{참고자료 목록}
\label{tab:references}
\begin{tabularx}{\textwidth}{c|l|l|X}
\toprule
\textbf{\#} & \textbf{자료명} & \textbf{유형} & \textbf{비고} \\
\midrule
1 & 시뮬레이터 v1.2.4.1 & HTML 파일 & 실행 파일 \\
2 & 개발일지 v1.2.4.1 & 문서 & 기술 상세 \\
3 & 퀵가이드 v1.2.4.1 & 문서 & 빠른 시작 안내 \\
4 & 문서버전관리가이드 v1.2.2 & 문서 & 버전 관리 체계 \\
\bottomrule
\end{tabularx}
\end{table}

%========================================
\section{생성/수정 이력}
%========================================

\begin{table}[H]
\centering
\caption{문서 버전 이력}
\label{tab:version_history}
\begin{tabularx}{\textwidth}{l|l|l|l|X}
\toprule
\textbf{버전} & \textbf{날짜} & \textbf{시간} & \textbf{변경 수준} & \textbf{변경 내용} \\
\midrule
v1.0.0.0 & 2026-01-31 & 16:00 & 최초 & 초안 작성 \\
v1.2.0.4 & 2026-02-01 & 15:00 & 마이너 & 읽기 순서 가이드, 빠른 시작 추가 \\
v1.2.1.0 & 2026-02-01 & 18:00 & 마이너 & 초보자 개선 (Tip 박스, FAQ, 용어집) \\
v1.2.1.1 & 2026-02-01 & 18:40 & 패치 & LaTeX 버전 변환, 그림/표 번호 추가 \\
\bottomrule
\end{tabularx}
\end{table}

\end{document}
